% Gerado por scripts/migrate_markdown_to_latex.py.
% Fonte migrada: fases/01_validacao_conceitual/docs/perguntas_pesquisa.md


\chapter{Perguntas de pesquisa e criterios de validade}




Este documento formaliza o que o estudo pretende responder e quais evidencias sao necessarias para uma conclusao ser aceita no artigo.



\section{Perguntas de pesquisa}



\begin{longtable}{@{}>{\RaggedRight\arraybackslash}p{0.205\linewidth}>{\RaggedRight\arraybackslash}p{0.205\linewidth}>{\RaggedRight\arraybackslash}p{0.205\linewidth}>{\RaggedRight\arraybackslash}p{0.205\linewidth}@{}}
\toprule
ID & Pergunta & Metricas associadas & Evidencia minima \\
\midrule
\endfirsthead
\toprule
ID & Pergunta & Metricas associadas & Evidencia minima \\
\midrule
\endhead
RQ1 & Quantizacao int8 preserva melhor a saida do bloco Transformer simplificado do que pruning por magnitude? & MSE, MAE, R2, similaridade de cosseno & Saida baseline, saida candidata e metricas de erro. \\
RQ2 & Quantizacao e pruning reduzem memoria teorica ou real no ambiente do estudo? & \texttt{max\_\allowbreak{}memory\_\allowbreak{}bytes} & Ambiente registrado, medicao ou estimativa de memoria e nivel de validade. \\
RQ3 & Alguma tecnica reduz latencia ou aumenta throughput sem degradacao excessiva da saida? & Latencia media, p50, p95, throughput, MSE e cosseno & Warmup, repeticoes, sincronizacao em CUDA e baseline comparavel. \\
RQ4 & Quando uma tecnica e apenas conceitual/numerica e quando pode ser chamada de hardware? & \texttt{validity\_\allowbreak{}level} & Matriz de validacao, dtype/formato e kernel/caminho de execucao. \\
\bottomrule
\end{longtable}



\section{Hipoteses}



\begin{landscape}
\scriptsize
\begin{longtable}{@{}>{\RaggedRight\arraybackslash}p{0.164\linewidth}>{\RaggedRight\arraybackslash}p{0.164\linewidth}>{\RaggedRight\arraybackslash}p{0.164\linewidth}>{\RaggedRight\arraybackslash}p{0.164\linewidth}>{\RaggedRight\arraybackslash}p{0.164\linewidth}@{}}
\toprule
ID & Pergunta & Hipotese & Confirma se & Rejeita se \\
\midrule
\endfirsthead
\toprule
ID & Pergunta & Hipotese & Confirma se & Rejeita se \\
\midrule
\endhead
H1 & RQ1 & Quantizacao int8 tera menor erro numerico do que pruning por magnitude no bloco controlado. & MSE e MAE menores, R2 e cosseno maiores que pruning na mesma configuracao. & Pruning empata ou supera quantizacao nas metricas de fidelidade. \\
H2 & RQ2 & Quantizacao int8 reduz memoria representacional; pruning por mascara densa nao prova reducao real de hardware. & Quantizacao apresenta menor armazenamento observado/estimado e pruning sem sparse kernel fica sem nivel hardware. & Quantizacao nao reduz armazenamento ou pruning demonstra formato/kernel esparso exploravel. \\
H3 & RQ3 & Ganhos de latencia/throughput dependem do caminho real de execucao, nao apenas da tecnica declarada. & Cenarios sem kernel/formato compativel nao sao promovidos a hardware mesmo que tenham erro aceitavel. & Ha evidencia de kernel/formato otimizado e medicao consistente de ganho fisico. \\
H4 & RQ4 & A matriz de validacao separa corretamente conclusoes conceituais, numericas, algoritmicas e de hardware. & Toda conclusao aponta para registro CSV e linha conceitual da matriz. & Alguma conclusao exige evidencia que nao esta registrada. \\
\bottomrule
\end{longtable}
\end{landscape}



\section{Regra de conclusao}



Uma conclusao so pode entrar no texto final se ela apontar para:


\begin{enumerate}
\item uma pergunta de pesquisa;
\item uma hipotese ou criterio de verificacao;
\item uma linha da matriz de validacao;
\item um resultado CSV ou teste deterministico;
\item um nivel de validade coerente com a evidencia observada.
\end{enumerate}



Essa regra e implementada no codigo por \texttt{validacao.\allowbreak{}pesquisa}, \texttt{validacao.\allowbreak{}benchmark\_\allowbreak{}controlado} e \texttt{validacao.\allowbreak{}analise\_\allowbreak{}resultados}.
